% Final report skeleton -- Cybersecurity Lab 1, ELTE.
% Build: make report   (runs analysis/make_figures.py first, then latexmk)
% Tables in tables/ and figures in ../figures/ are generated from results/raw/;
% never edit them by hand -- re-run the analysis instead.
\documentclass[11pt,a4paper]{article}
\usepackage[margin=2.5cm]{geometry}
\usepackage{microtype,booktabs,graphicx,amsmath,hyperref,xcolor}
\usepackage[capitalise]{cleveref}
\graphicspath{{../figures/}}
\newcommand{\todo}[1]{\textcolor{red!70!black}{[TODO: #1]}}
% Include a generated table if it exists, otherwise a placeholder.
\newcommand{\gentable}[1]{\IfFileExists{tables/#1.tex}{\input{tables/#1.tex}}{\todo{run the experiment to generate \detokenize{#1}}}}
\newcommand{\genfigure}[3]{\IfFileExists{../figures/#1.pdf}{%
  \begin{figure}[htbp]\centering\includegraphics[width=#2\linewidth]{#1}\caption{#3}\label{fig:#1}\end{figure}}%
  {\todo{run the experiment to generate figure \detokenize{#1}}}}

\title{A Comparative Study of Post-Quantum Key-Encapsulation Mechanisms\\and Their Implementations}
\author{Hamdan Al-Jabri\\ELTE Faculty of Informatics, Cybersecurity Lab 1}
\date{Autumn 2026}

\begin{document}
\maketitle

\begin{abstract}
\todo{150 words: the problem, what was compared, the four headline results (one per RQ), the takeaway.}
\end{abstract}

\section{Introduction}\label{sec:intro}
\todo{Harvest-now-decrypt-later; NIST standardisation (FIPS 203 in Aug 2024, HQC selected Mar 2025); why
implementation matters as much as the algorithm.}

Research questions:
\begin{description}
  \item[RQ1 Cost] How do keygen, encapsulation and decapsulation compare in time, object size and stack use
    across families and security levels, and does implementation optimisation change the ranking?
  \item[RQ2 Deployability] What does each KEM add to a TLS 1.3 handshake in latency and bytes on the wire
    under realistic delay and loss?
  \item[RQ3 Implementation security] Do the tested builds show secret-dependent timing, and can a known flaw
    (KyberSlash) be detected with the same tooling?
  \item[RQ4 Correctness] Do independent implementations agree on official test vectors and interoperate?
\end{description}
\todo{Contributions paragraph.}

\section{Background}\label{sec:background}
\subsection{KEMs and IND-CCA2 security}
\todo{Definition; Fujisaki--Okamoto transform; implicit rejection.}
\subsection{The four families}
\todo{ML-KEM (Module-LWE), FrodoKEM (plain LWE), HQC (quasi-cyclic codes), Classic McEliece (binary Goppa
codes); hardness assumptions; decryption-failure rates; standardisation status.}
\subsection{Timing side channels}
\todo{Constant-time programming; compiler-introduced leaks; KyberSlash (secret-dependent division).}

\section{Related work}\label{sec:related}
\todo{Earlier benchmark studies (eBACS/SUPERCOP, pqm4, liboqs profiling); TLS experiments by Cloudflare and
Google; dudect, ctgrind/TIMECOP, KyberSlash.}

\section{Methodology}\label{sec:method}
\subsection{Implementations and platforms}
\todo{liboqs 0.16.0 ref (portable C) and opt (AVX2 on x86-64, NEON on ARM64) builds; OpenSSL 3.5.8;
oqs-provider 0.12.0-rc2; the two machines (\texttt{make sysinfo}); pinned versions in \texttt{versions.env}.}
\subsection{E1 microbenchmarks}
\todo{Per-call timing with the portable timer (RDTSC / CNTVCT\_EL0), median and IQR, warm-up, CPU pinning,
cool-down on the fanless laptop, stack painting.}
\subsection{E2 TLS handshakes}
\todo{Rootless network namespace, loopback MTU 1500, \texttt{tc netem} delay and loss, one group offered per
handshake, ECDSA P-256 certificate, bytes from \texttt{tcp\_info}.}
\subsection{E3 constant-time analysis}
\todo{dudect fixed-vs-random decapsulation with planted positive/negative controls; KyberSlash control at four
optimisation levels plus a static DIV scan; Valgrind with random bytes marked secret, compared against liboqs'
known-issue list. Note the dead-code pitfall found and fixed in the harness.}
\subsection{E4 correctness}
\todo{liboqs KATs, NIST ACVP FIPS 203 vectors, cross-build and liboqs$\leftrightarrow$OpenSSL interop.}
\subsection{Threats to validity}
\todo{Single machine per architecture; VM timer resolution on ARM (41.7\,ns ticks); Turbo Boost / thermal
throttling; dudect ``no evidence'' is not proof; netem on loopback is not a real network.}

\section{Results}\label{sec:results}
\subsection{RQ1: cost}
\genfigure{fig_e1_ops}{1.0}{Median time per operation (log scale).}
\gentable{tab_sizes}
\todo{Interpret: ranking, effect of optimisation per family, OpenSSL vs liboqs ML-KEM, x86 vs ARM.}

\subsection{RQ2: deployability in TLS}
\genfigure{fig_e2_handshake}{1.0}{Median TLS 1.3 handshake time per group.}
\gentable{tab_e2_bytes}
\todo{Interpret: bytes vs packets vs latency; loss sensitivity; why Classic McEliece has no TLS group
(key share limited to 65{,}535 bytes).}

\subsection{RQ3: implementation security}
\genfigure{fig_e3_dudect}{0.8}{dudect maximum $|t|$ per target; 10 is the leak threshold.}
\gentable{tab_e3_valgrind_x86_64}
\gentable{tab_e3_divscan_x86_64}
\gentable{tab_e3_divscan_aarch64}
\todo{Interpret: controls first; KEM verdicts; KyberSlash compiler/CPU dependence; Valgrind vs upstream issues.}

\subsection{RQ4: correctness}
\gentable{tab_e4_x86_64}
\gentable{tab_e4_aarch64}
\todo{All checks passed / any failures.}

\section{Discussion}\label{sec:discussion}
\todo{Which KEM fits which use case (TLS, VPN, embedded, long-term static keys); hybrid vs pure; surprises.}

\section{Limitations and future work}\label{sec:limits}
\todo{Power/EM side channels, fault injection, embedded targets (pqm4), more CPUs.}

\section{Conclusion}\label{sec:conclusion}
\todo{One answer per RQ.}

\appendix
\section{Reproducing the results}
All numbers come from \url{https://github.com/<your-account>/pq-kem-study} at the commit recorded in each CSV
header: \texttt{make deps \&\& make all}.

\end{document}
